\paragraph{Simultaneous recoloring from capture probes.}
An input can be irrelevant after one fixing but live at both endpoints
of a simultaneous recoloring. The following profile argument lets us
use such intermediate capture probes without choosing a path through them.
For a valuation $w$, write $E_s(w)=\{i:w_i=\{s\}\}$.

\begin{lemma}[Marked profile transfer]
\label{core:lem:marked-profile-transfer}
Let $H$ be core, with $\Phi_r(H)(A)=1$ and $\Phi_r(H)(B)=0$.
Fix an observed winner $k$. For each $j\in A\setminus B$, suppose
there is a valuation $w^j$ with
\[
 k\in w^j_j,\qquad E_r(w^j)\subseteq B,\qquad
 E_s(w^j)\cap A=\varnothing\quad(s\ne r).
\]
Then $k\in H(w^j)$ for at least one such $j$.
\end{lemma}
\begin{proof}
Suppose otherwise. Descend a word for $H$, preserving its positive
$r$-profile at $A$, its negative $r$-profile at $B$, and omission of
$k$ on all the displayed rows.
At an $r$-root choose a child positive at $A$; both children are
negative at $B$. Neither can be an $r$-absorber on a row whose
$r$-coalition lies in $B$, so rowwise omission passes to both.
At a foreign $s$-root both children are positive at $A$; choose one
negative at $B$. Neither child can be singleton $s$ on a row:
its winning $s$-coalition would be disjoint from its winning
$r$-coalition $A$, contrary to \cref{core:thm:profiles}.
Thus omission again passes to both children.
A resulting leaf must be indexed by $j\in A\setminus B$ and
contains $k$ on its own designated row, a contradiction.
\end{proof}

\begin{theorem}[A simultaneous capture-probe rule]
\label{core:thm:simultaneous-capture-probes}
Let $K\geq3$, let $C$ be core on $V$, and let $p,q$ be proper
colorings. Put $S=\{j:p_j\ne q_j\}$,
where a status is an owner or the uncolored status.
Suppose that for each $j\in S$ there is a coloring $e^j$ of
$V\setminus\{j\}$ with $P_j^C(e^j)=0$ and
\[
 e^j_r\subseteq
 \begin{cases}
 q_r,&r=p_j,\\
 p_r,&r=q_j,\\
 p_r\cap q_r,&r\notin\{p_j,q_j\}.
 \end{cases}
\]
Then $p\sim_Cq$.
\end{theorem}
\begin{proof}
Choose an owner $k$ missing at $q$ and put $N=I\setminus\{k\}$.
Every $e^j$ uses only the intersection of the endpoint palettes,
so also omits $k$. Form the negative row $R_j$ by giving $X$ menu
$N$, input $j$ menu $I$, and all other inputs their $e^j$ atom or
$N$ when uncolored. The capture assumption makes $C$ omit $k$ there.
Call a word quiet if it omits $k$ on every $R_j$.

Fix any common assignment $v$ to the passive inputs and any menu
$M$ for $X$. At one positive evaluation overwrite the colored slots
of $v$ by $p$; at the other overwrite them by $q$, using $X=M$
in both. Write $A_r,B_r$ for their atomic coalitions, including the
common $X$ flag if $M=\{r\}$. Then $p_r\subseteq A_r$ and
$q_r\subseteq B_r$.
Every $j\in A_r\setminus B_r$ lies in $S$: unchanged inputs have
identical menus. Its original $p$-status is either $r$ or uncolored,
and its original $q$-status is not $r$. The capture bounds imply
\[
 e^j_r\subseteq q_r\subseteq B_r,\qquad
 e^j_s\subseteq p_s\quad(s\ne r).
\]
The latter sets are disjoint from $A_r$, since their inputs remain
fixed at $s$. By \cref{core:lem:marked-profile-transfer}, every quiet
word therefore has identical singleton flags at the positive
evaluations. Interchanging the endpoints proves the reverse transfer.

Induct on a word to show that quietness forces equality of these
positive outputs. The distinguished leaf has the same menu $M$ at
both. An unchanged passive leaf also has identical menus; a changed
leaf $j$ is not quiet on $R_j$. At a root, atomic positive outputs
agree by profile transfer. Otherwise both are ambiguous, so neither
child is the root owner's singleton at either positive evaluation.
Each negative atomic coalition $e^j_r$ lies in $p_r$ or $q_r$,
hence in $A_r$ or $B_r$, ruling out that absorber there too.
Quietness passes to both children; induction gives equal child
values, and the root takes union at both positive evaluations.

Since $v$ and $M$ were arbitrary, every quiet word has equal full
clamped cofactors under $p,q$. Any $F\in\mathcal I(C)$ is quiet
by its upper bound $C+X$. Its two cofactors are therefore either
both independent of $X$ or both dependent on it. By
\cref{core:cor:colored-cofactor-dependence}, this says exactly
$P_F(p)=0$ if and only if $P_F(q)=0$, proving $p\sim_Cq$.
\end{proof}

\paragraph{An exact kernel for graph sums.}
Let $G=(V,E)$ be a finite simple graph with at least one edge.
Assign an arbitrary owner $r_{ij}$ to each edge, and put
\[
 D(v)=\sum_{\{i,j\}\in E}v_i\star_{r_{ij}}v_j.
\]
Isolated vertices are allowed and are ignored by $D$.

\begin{theorem}[The activation kernel of a graph sum]
\label{core:thm:graph-activation-kernel}
For $K\geq3$, proper partial colorings $p,q$ are activation-equivalent
above $D$ exactly when every edge $\{i,j\}$ satisfies
\[
 r_{ij}\in\{p_i,p_j\}\cap\{q_i,q_j\}
 \quad\text{or}\quad p_i=q_i\ \text{and}\ p_j=q_j.
\]
\end{theorem}
\begin{proof}
Each edge summand is a core minorant of $D$, so its interval
predicates lift by adjoining $D$. The kernel for a single owned
choice, given by \cref{core:thm:profile-source-interval}, has one
class where some endpoint is its owner, and singleton classes on
its two endpoint statuses elsewhere. This proves necessity.

For sufficiency, at every changed vertex $j$ color each neighbor $i$ by its edge
owner $r_{ij}$ and leave all other inputs uncolored. Simplicity
makes this a well-defined coloring $e^j$. Every incident edge then
absorbs the source $j$, so $P_j^D(e^j)=0$.
The local edge test gives exactly the required inclusions:
\[
 e^j_r\subseteq
 \begin{cases}
 q_r,&r=p_j,\\
 p_r,&r=q_j,\\
 p_r\cap q_r,&r\notin\{p_j,q_j\}.
 \end{cases}
\]
For example, an $r$-neighbor of a vertex leaving $r$ must enter
$r$, and an $r$-neighbor of a vertex entering $r$ must leave it.
If neither status is $r$, the neighbor must be unchanged at $r$.
The bounds also make each $e^j$ proper, since its palette lies in
both proper endpoint palettes. The simultaneous capture-probe rule
therefore proves equivalence.
\end{proof}

\begin{proposition}[Graph comparisons need not factor edgewise]
\label{core:prop:graph-order-nonfactorization}
For distinct owners $r,s$, put
\[
 D=(x\star_r y)+(y\star_r z),\qquad
 p=(s,\circ,\circ),\quad q=(r,\circ,\circ),
\]
where $\circ$ is uncolored. The restrictions satisfy
$p\preceq q$ above each edge summand, but $p\not\preceq_Dq$.
\end{proposition}
\begin{proof}
On the first edge, $(s,\circ)$ extends to $(s,r)$, which is
equivalent to $(r,\circ)$ in the owning singleton class.
Antitonicity gives the first local comparison; the second edge has
identical restrictions. For the global separation take
\[
 S=((X\star_s x)\star_r y)\star_r z,\qquad F=D+S.
\]
If $y$ or $z$ is the $r$ atom, $S=\{r\}\subseteq D$.
Otherwise the second edge is $y+z$ and the first contains $x$,
so $D$ is the full input union. Thus $S\subseteq D+X$ in every
case and $F\in\mathcal I(D)$. Its distinguished predicate is
\[
 [x\text{ is not }s]\,[y\text{ is not }r]\,[z\text{ is not }r],
\]
which is one at $q$ and zero at $p$.
\end{proof}

For distinct owners $a,b$, if all edge labels and input colors belong
to $\{a,b\}$, each edge
is Boolean: an $a$-edge takes OR and a $b$-edge takes AND, with
$a=1,b=0$. Here the whole activation preorder is equality of edge-value
vectors. Indeed $D+(X\star_a H)$ and $D+(X\star_b H)$, for an
edge summand $H$, separate unequal Boolean values in both directions.
Equal vectors satisfy the theorem's local test.
Each class is a monotone 2-CNF solution set: an $a$-edge with output
$1$ gives $x_i\lor x_j$, and output $0$ fixes both endpoints to $0$;
a $b$-edge with output $0$ gives $\neg x_i\lor\neg x_j$, and
output $1$ fixes both endpoints to $1$.

The simple-graph hypothesis matters. Oppositely labelled parallel
edges on one pair give $D=x+y$. The swapped colorings $(a,b)$ and
$(b,a)$ agree in both local owner classes, yet
$D+(X\star_a x)$ separates them.

\begin{theorem}[Activation below an owning coloring]
\label{core:thm:bipartite-owning-activation}
Suppose $G$ is bipartite and all its edge labels belong to two
distinct owners $a,b$. Let $q$ be any proper coloring with an
owner-colored endpoint at every edge. Then for every proper $p$,
\[
 p\preceq_Dq\quad\Longleftrightarrow\quad
 \forall\{i,j\}\in E,\ \forall s\ne r_{ij},\quad
 \neg(p_i=p_j=s).
\]
Thus no edge may be monochromatic in a color different from its owner.
\end{theorem}
\begin{proof}
By the graph-kernel theorem, replace $q$ by the coloring giving
$a$ to one bipartition class $L$ and $b$ to the other class $R$.
The swapped coloring $q'$ is equivalent too.
Let $T=\Theta_{D,q}$, realized by a primitive one-$X$ spine
$S\subseteq D+X$. Choose $c\notin\{a,b\}$ and put
$N=I\setminus\{c\}$. At row $R_j$ give input $j$ menu $I$,
each neighbor its incident edge owner, and all other inputs and
$X$ menu $N$. Every incident edge absorbs the marked source,
so $D+X$ and hence $S$ omit $c$ there.

Every $a$- or $b$-owned path helper avoids its owner's singleton
on both $L$ and $R$, since $T$ is positive at $q,q'$.
Its owner's atomic coalition at $R_j$ lies in one bipartition
class, so it cannot absorb there either. Other owners have no
atomic inputs on $R_j$. Starting at ambiguous $X=N$, all path
gates therefore take union on these rows. Each helper must itself
omit $c$ on every $R_j$.

Fix an owner $s$ and a set $A$ containing no edge with owner
different from $s$. No $s$-owned helper can have
$\Phi_s(H)(A)=1$. For $s=a$, it has negative profile at $L$.
Rows indexed by $j\in A\setminus L$ have $a$-coalition inside
$L$ and $b$-coalition disjoint from $A$, since $A$ contains no
$b$-edge. Marked profile transfer forces $c$ in a helper output,
a contradiction. The case $s=b$ is symmetric, using $R$.
For other $s$, take the negative coalition empty: $A$ is independent
in the whole graph, so every neighbor coalition at a row marked
in $A$ is disjoint from $A$. The same transfer again contradicts
omission of $c$.
Thus $T_s(A)=1$ for all such $A$.

Conversely, if $A$ contains an edge $e$ owned by $r\ne s$,
its summand $H_e$ is a minorant of $D$.
The interval $D+(X\star_s H_e)$ is positive at $q$ but zero on
the monochromatic $s$-coloring of $A$. Hence $T_s(A)=0$.
The resulting exact factors give the stated equivalence.
\end{proof}

Bipartiteness cannot be omitted. On the triangle
$D=(x\star_a y)+(y\star_a z)+(z\star_b x)$, the target
$q=(b,a,b)$ owns every edge, while $t=(\circ,c,\circ)$ has no
monochromatic edge. Yet the interval
\[
 D+\bigl((X\star_c y)\star_a(x\star_b z)\bigr)
\]
has distinguished predicate positive at $q$ and zero at $t$.
To check its upper bound, if
$x=z=a$ its last gate returns $a$; if $y=a$, both passive helper
contributions belong to $D$; otherwise $D$ is the full input union.

\begin{proposition}[Nonlocal activation and ownerwise mixing]
\label{core:prop:activation-kernel-cycle}
For every $K\geq3$, there is a four-input core operation with two
activation-equivalent proper colorings, each isolated in the graph
of capture edges. For distinct owners $a,b$, take
\[
 D=(x\star_a y)+(y\star_b z)+(z\star_a w)+(w\star_b x),
 \qquad p=(a,b,a,b),\quad q=(b,a,b,a).
\]
For $c\notin\{a,b\}$, the monochromatic point
$t=(c,\circ,\circ,\circ)$ lies strictly below both $p,q$.
Yet no coloring, even one using every owner, has $x=c$ and
an owner-colored endpoint at every edge.
\end{proposition}
\begin{proof}
Every edge has its owner at one endpoint under both colorings, so
\cref{core:thm:graph-activation-kernel} gives $p\sim_Dq$.
Every vertex has an incident edge of each label. Under either
alternating coloring one incident edge leaves that vertex's source
live. Both colorings are fully colored, so neither has a capture edge.

The square is bipartite, and $t$ has no monochromatic edge.
Thus \cref{core:thm:bipartite-owning-activation} gives
$t\preceq_Dp$.
The interval $D+(X\star_a(x\star_a y))$ has predicate one at
$t$ and zero at $p$, excluding the reverse comparison.
Finally $x=c$ would successively force $y=a,z=b,w=a$ to cover
the first three edges, leaving the final $b$-edge uncovered.
\end{proof}
